\subsubsection{Розробка тестового застосунку з використанням компонентів
Loki23}\label{ux440ux43eux437ux440ux43eux431ux43aux430-ux442ux435ux441ux442ux43eux432ux43eux433ux43e-ux437ux430ux441ux442ux43eux441ux443ux43dux43aux443-ux437-ux432ux438ux43aux43eux440ux438ux441ux442ux430ux43dux43dux44fux43c-ux43aux43eux43cux43fux43eux43dux435ux43dux442ux456ux432-loki23}

Для оцінювання бібліотеки Loki23 у прикладному сценарії було створено
тестовий застосунок \texttt{netprobe}. Це консольна програма для
базового аналізу мережевого трафіку, яка може отримувати пакети з
мережевого інтерфейсу або з попередньо записаного \texttt{pcap}-файлу.
Основне призначення застосунку полягає не у створенні повноцінної
системи мережевого моніторингу, а у перевірці того, як компоненти Loki23
поводяться в реальній структурі C++23-програми, де є динамічні ресурси,
поліморфізм, потоки виконання, фабричне створення об'єктів і
диспетчеризація за фактичними типами.

Вибір мережевого аналізатора як тестового прикладу є обґрунтованим з
кількох причин. По-перше, мережевий трафік природно складається з
неоднорідних об'єктів: TCP-пакети, UDP-пакети, DNS-запити,
DNS-відповіді, HTTP-запити та HTTP-відповіді мають спільні властивості,
але потребують різної логіки обробки. По-друге, захоплення трафіку
пов'язане з ресурсами, які необхідно відкривати та закривати у
правильному порядку. По-третє, отримання пакетів і їх подальший аналіз
зручно розділяти на незалежні компоненти. Усі ці особливості добре
відповідають задачам, для яких у Loki23 реалізовано політично керовані
розумні вказівники, списки типів, фабрики і мультиметоди.

У загальному вигляді \texttt{netprobe} складається з трьох частин. Перша
частина відповідає за отримання сирих пакетів. Для цього
використовується бібліотека PcapPlusPlus, яка надає доступ до мережевих
інтерфейсів і файлів захоплення. Друга частина перетворює сирі дані на
об'єктну модель застосунку. Замість того щоб працювати з низькорівневими
структурами PcapPlusPlus у всіх компонентах програми, \texttt{netprobe}
створює власні класи пакетів із потрібними для аналізу полями. Третя
частина складається зі слухачів, які отримують уже розпізнані пакети і
виконують окремі види обробки: збір статистики, виведення в консоль або
пошук простих зв'язків між DNS-відповідями та подальшими з'єднаннями.

Центральним елементом моделі даних є абстрактний тип \texttt{NetPacket}.
Він задає спільний інтерфейс для всіх пакетів, що обробляються
застосунком. У ньому зберігаються IP-адреси джерела і призначення,
розмір пакета та час отримання. Крім базових полів, клас визначає
віртуальні операції для отримання назви типу, створення поліморфної
копії та формування текстового подання. Наявність методу поліморфного
копіювання є важливою деталлю: пакети передаються через базовий тип, але
при цьому слухачі повинні отримувати об'єкт зі збереженим конкретним
типом, наприклад DNS або HTTP.

Від \texttt{NetPacket} успадковуються класи, що відповідають окремим
протоколам. \texttt{TcpPacket} містить порти, номери послідовності й
підтвердження, розмір вікна та набір прапорців TCP. \texttt{UdpPacket}
зберігає порти UDP. \texttt{DnsPacket} описує DNS-повідомлення:
ідентифікатор запиту, домен, тип запису, ознаку відповіді та список
отриманих адрес. \texttt{HttpPacket} є спеціалізованим TCP-пакетом і
додає дані прикладного рівня: HTTP-метод, URL, версію протоколу,
заголовок \texttt{Host} або код статусу відповіді. Така ієрархія
дозволяє обробляти пакети узагальнено, але за потреби звертатися до
протокольних деталей.

Отримання пакетів винесено в інтерфейс \texttt{PacketCapturer}. Він не
прив'язаний до конкретного джерела трафіку і визначає лише базові
операції відкриття, закриття та передавання пакета підписаним слухачам.
У застосунку реалізовано два класи, що відповідають цьому інтерфейсу.
\texttt{LiveCapturer} працює з активним мережевим інтерфейсом і запускає
асинхронне захоплення пакетів через PcapPlusPlus.
\texttt{PcapFileCapturer} читає пакети з файлу. Завдяки такому
розділенню решта програми не залежить від того, звідки саме надходить
трафік.

Під час розбору сирого пакета застосунок послідовно намагається створити
один із відомих типів: DNS, HTTP, TCP або UDP. Порядок перевірки має
практичне значення. HTTP-пакет одночасно є TCP-пакетом, тому його
потрібно розпізнавати до загального TCP-випадку. Інакше інформація
прикладного рівня була б втрачена, а пакет потрапив би лише до загальної
TCP-статистики. Такий підхід є простим, але достатнім для тестового
застосунку, оскільки набір підтримуваних протоколів є обмеженим і явно
заданим у коді.

Перший компонент Loki23, який активно використовується в
\texttt{netprobe}, - це \texttt{SmartPtr}. Для різних об'єктів
застосунок використовує різні політики володіння. Об'єкти захоплення
пакетів і слухачі створюються один раз, мають єдиного власника і не
повинні копіюватися. Для них використовується політика \texttt{NoCopy}.
Вона робить намір володіння явним і забороняє випадкове копіювання на
рівні типів. Це важливо для компонентів, які пов'язані з потоками або
зовнішніми ресурсами, оскільки копіювання такого об'єкта могло б
спричинити дублювання стану або некоректне завершення роботи.

Інша ситуація виникає під час передавання пакетів слухачам. Один
розпізнаний пакет може бути потрібний кільком компонентам одночасно:
статистичному слухачу, розумному слухачу або журналюванню. При цьому
кожен слухач працює у власному потоці і може завершити обробку в різний
момент. Для такого сценарію використовується
\texttt{SmartPtr\textless{}NetPacket,\ RefCountedMT\textgreater{}}.
Потокобезпечний лічильник посилань дозволяє передати одну поліморфну
копію пакета кільком отримувачам і звільнити пам'ять лише тоді, коли
останній із них завершить роботу з об'єктом.

У цьому прикладі \texttt{SmartPtr} демонструє основну ідею політично
керованого проєктування. Поведінка вказівника не є фіксованою для всіх
випадків. Вона визначається політикою, яка передається як параметр
шаблону. Завдяки цьому один і той самий інструмент може описувати різні
моделі володіння: унікальне некопійоване володіння для компонентів
застосунку і спільне потокобезпечне володіння для пакетів, що
розсилаються між слухачами.

Другим важливим механізмом Loki23 у застосунку є фабричне створення
об'єктів на основі списків типів. У \texttt{main.cpp} визначено два
набори типів: список слухачів і список джерел пакетів. До першого
входять \texttt{PacketLogger}, \texttt{StatisticsListener} і
\texttt{SmartListener}, до другого - \texttt{LiveCapturer} і
\texttt{PcapFileCapturer}. Ці набори задано через \texttt{type\_list},
тобто як структури, відомі на етапі компіляції. Далі вони
використовуються для побудови \texttt{HierarchyFactory}.

Фабрика дозволяє створювати конкретні об'єкти без розміщення прямого
\texttt{new} в основній логіці програми. Наприклад, статистичний слухач
створюється через поле фабрики, пов'язане з типом
\texttt{StatisticsListener}, після чого зберігається в
\texttt{SmartPtr\textless{}PacketListener,\ NoCopy\textgreater{}}. Для
основної функції важливо не те, як саме виділяється пам'ять, а те, що
отриманий об'єкт реалізує інтерфейс \texttt{PacketListener}. Такий
підхід зменшує зв'язаність між кодом запуску програми і конкретними
класами компонентів.

Використання \texttt{type\_list} у цьому місці має ще одну перевагу:
набір допустимих реалізацій зафіксовано явно. Якщо потрібно додати новий
тип слухача, наприклад аналізатор підозрілих TCP-з'єднань, потрібно
реалізувати клас із відповідним інтерфейсом і додати його до списку
типів. Подібно можна додати нове джерело пакетів, не змінюючи логіку
наявних слухачів. Отже, фабрика і списки типів у \texttt{netprobe}
працюють як засіб розширення архітектури без дублювання коду створення
об'єктів.

Обробка пакетів побудована за моделлю підписників. Клас
\texttt{PacketCapturer} зберігає список вказівників на
\texttt{PacketListener} і після розпізнавання пакета розсилає його всім
зареєстрованим слухачам. Сам слухач має внутрішню чергу, м'ютекс, умову
очікування і робочий потік. Коли пакет надходить у метод \texttt{push},
він додається до черги, а робочий потік пробуджується та викликає метод
\texttt{on\_packet}. Завдяки цьому конкретні слухачі не реалізують
синхронізацію самостійно, а описують лише власну прикладну поведінку.

Таке рішення також зменшує навантаження на callback захоплення пакетів.
У живому режимі PcapPlusPlus викликає функцію обробки при надходженні
пакета. Якщо в цій функції виконувати складний аналіз, захоплення може
затримуватися. У \texttt{netprobe} callback лише створює об'єкт пакета і
передає його слухачам, а основна робота виконується окремо. Для
тестування Loki23 це корисно, оскільки дозволяє перевірити не тільки
базове володіння пам'яттю, а й поведінку об'єктів у багатопотоковому
сценарії.

Найпростіший слухач, \texttt{PacketLogger}, призначений для
налагодження. Він виводить текстове представлення кожного пакета. Його
можна підключити під час розробки, щоб переконатися, що розбір пакетів
працює очікувано і що об'єкти моделі заповнюються правильними
значеннями.

\texttt{StatisticsListener} виконує агрегований аналіз. Для TCP він
накопичує кількість пакетів, обсяг даних і кількість пакетів із
прапорцями \texttt{SYN}, \texttt{FIN} та \texttt{RST}. Для UDP
зберігаються кількість пакетів і байтів. Для DNS фіксуються запити,
відповіді, типи записів і доменні імена. Для HTTP окремо враховуються
запити, відповіді, методи та коди статусу. Після завершення захоплення
слухач формує підсумковий звіт. Цей компонент демонструє типовий
сценарій моніторингу, у якому окремі пакети не зберігаються довго, а
використовуються для оновлення статистичних показників.

\texttt{SmartListener} реалізує складнішу поведінку і демонструє
мультиметоди Loki23. Він використовує \texttt{DynamicDispatcher}, до
якого реєструються обробники для конкретних типів пакетів. Якщо
надходить DNS-відповідь, слухач запам'ятовує відповідність між отриманою
IP-адресою і доменом. Якщо надходить HTTP-запит або TCP SYN-пакет,
слухач намагається зіставити адресу призначення з уже відомим доменом і
виводить повідомлення про ймовірне відвідування цього ресурсу. У
прикладному коді при цьому не потрібно писати довгий ланцюжок перевірок
типу; вибір обробника делеговано диспетчеру.

Крім динамічної диспетчеризації, \texttt{SmartListener} використовує
\texttt{StaticDispatcher} для аналізу пари послідовних пакетів. Слухач
зберігає попередній пакет, а при надходженні нового перевіряє комбінації
типів. Наприклад, якщо DNS-відповідь безпосередньо передує TCP, UDP або
HTTP-з'єднанню до однієї з адрес із відповіді, застосунок повідомляє про
такий зв'язок. Цей приклад показує різницю між двома підходами:
\texttt{DynamicDispatcher} зручний для реєстрації обробників окремих
об'єктів, а \texttt{StaticDispatcher} - для наперед відомих комбінацій
типів.

Основна функція програми виконує роль композиційного кореня. У ній
створюються фабрики, через них будуються слухачі та потрібний capturer,
після чого слухачі підписуються на отримання пакетів. Якщо користувач
передає параметр \texttt{-f}, програма читає трафік із файлу. Якщо
параметр не передано або використано \texttt{-i}, запускається живе
захоплення з мережевого інтерфейсу. У живому режимі програма працює
обмежений час, після чого закриває capturer і друкує статистичний звіт.

Режим роботи з файлом є найбільш придатним для повторюваної перевірки.
Один і той самий дамп можна аналізувати багато разів після змін у коді
Loki23, порівнюючи отримані результати. Живе захоплення, навпаки,
демонструє інтеграцію з реальною мережевою підсистемою, але його
результат залежить від поточного трафіку. Тому в контексті дипломної
роботи файл \texttt{pcap} є зручним засобом контрольного запуску, а
робота з інтерфейсом показує, що архітектура не обмежена лише
офлайн-аналізом.

Застосунок \texttt{netprobe} можна розглядати як інтеграційний тест для
Loki23. Він одночасно використовує кілька незалежних компонентів
бібліотеки і перевіряє їх у спільній роботі. \texttt{SmartPtr} керує
об'єктами з різними моделями володіння. \texttt{type\_list} описує
набори допустимих реалізацій. \texttt{HierarchyFactory} створює
компоненти за цими наборами. Диспетчери забезпечують вибір обробника на
основі фактичного типу пакета. Якщо будь-який із цих механізмів має
помилки в інтерфейсі або життєвому циклі, вони проявляються у достатньо
реалістичному прикладному коді.

Під час аналізу результатів роботи \texttt{netprobe} можна виділити
кілька критеріїв коректності. Перший критерій - правильне розпізнавання
протоколів і заповнення статистики. DNS, HTTP, TCP і UDP-пакети повинні
потрапляти у відповідні категорії, а специфічні поля мають враховуватися
лише для сумісних типів. Другий критерій - коректне завершення роботи:
після закриття capturer-об'єкта і знищення слухачів не повинно
залишатися активних потоків або незвільнених ресурсів. Третій критерій -
безпечна передача пакетів між потоками. Один пакет може оброблятися
кількома слухачами, тому механізм підрахунку посилань має забезпечити
правильний час життя об'єкта.

Разом із перевагами потрібно враховувати й обмеження реалізації.
\texttt{netprobe} не реконструює TCP-потоки, не аналізує зашифрований
HTTPS-трафік, не реалізує повноцінні правила виявлення атак і не
зберігає історію всіх мережевих подій. Кореляція DNS і подальших
з'єднань є спрощеною евристикою, а не точним доказом відвідування
певного ресурсу. Однак ці обмеження відповідають меті застосунку. Він
має бути достатньо реалістичним, щоб перевіряти бібліотеку в умовах,
наближених до практичного програмування, але не перетворюватися на
окрему велику систему мережевої безпеки.

Отже, \texttt{netprobe} демонструє практичне використання Loki23 у
модульному C++23-застосунку. На його прикладі видно, як шаблонні
компоненти можуть допомагати у створенні гнучкої архітектури: ресурси
мають явно описане володіння, набір реалізацій задається на етапі
компіляції, об'єкти створюються через фабрики, а обробка поліморфних
пакетів виконується через диспетчери. У межах дипломної роботи цей
застосунок підтверджує, що Loki23 може використовуватися не лише в
демонстраційних фрагментах коду, а й у цілісному прикладному сценарії з
потоками, зовнішньою бібліотекою і реальною предметною областю.
